\subsection{El operador límite y las dos hojas de paridad}\label{sec:D03}
Se usan la comparación uniforme con el perfil continuo (Sección~\ref{sec:P06}) y la brecha de la Sección~\ref{sec:P03}.
Sean $v_n=(\mathbf1_R,-\mathbf1_C,0_D,0_A)$ y
\[
 \begin{gathered}
 P_n=\frac{v_nv_n^T}{2n},\qquad J_n=S_n-I+P_n,\\
 S_n=\operatorname{diag}(0_R,0_C,h_D,h_A)+A^T\operatorname{diag}(g)A,
 \qquad g_c=\frac{\pi_c(1-\pi_c)}{1-4\pi_c}.
 \end{gathered}                                         \tag{D3.1}
\]
La matriz $S_n+P_n$ reemplaza por uno el único autovalor cero de $S_n$.
Así $\det(I+J_n)=\operatorname{pdet}S_n$ y
$I+J_n\ge\min(1,\gamma)I$ uniformemente.

\paragraph{Diagonal y traza.}
En todas las familias, $h_L+\sum_{c\in L}\pi_c=1$, con $h_R=h_C=0$.
Por tanto $(S_n)_{LL}=1+\delta_L$, con
$\delta_L=\sum_{c\in L}(g_c-\pi_c)=O(n^{-1})$.
La parte diagonal satisface $\sum_L\delta_L^2=O(n^{-1})$, pero
\[
 \operatorname{tr}J_n
 =1+4\sum_c(g_c-\pi_c)
 =1+12\sum_c\frac{\pi_c^2}{1-4\pi_c}
 \longrightarrow1+12I_2.                          \tag{D3.2}
\]
Aquí se usan $\sum\pi_c^3=O(n^{-1})$ y $\sum\pi_c^2\to I_2$.
La traza no se obtiene tomando la traza de un límite Hilbert--Schmidt.

\paragraph{Espacio fijo e isometrías.}
Con los intervalos $I_F$ de la Sección~\ref{sec:preliminares}, sea
\[
 \mathcal H=\bigoplus_{F=R,C,D,A}
          L^2(I_F\times\{0,1\},dt/2).
\]
Cada familia tiene dos hojas: la hoja de una línea es la paridad de su índice absoluto
$i,j,k,l$, respectivamente. En una casilla $(i,j)$ los índices diagonales $i-j$ e $i+j-1$
difieren en $2j-1$, de modo que siempre tienen paridades opuestas. Si se reuniesen las dos
paridades antes de construir el operador, se perdería esa información sobre qué diagonales
pueden cortarse en una casilla.
Para cada familia y paridad, con $m_{F,p}$ índices, se divide $I_F$
en $m_{F,p}$ intervalos iguales, las celdas, asignados a los índices por orden creciente.
La medida de cada celda en $\mathcal H$ es
$\mu_{F,p}=|I_F|/(2m_{F,p})=n^{-1}(1+O(n^{-1}))$.
La isometría desde el espacio euclídeo finito es
\[
 E_ne_L=\mu_{F,p}^{-1/2}\mathbf1_{\text{celda de }L}.
                                                        \tag{D3.3}
\]
Toda celda queda a distancia $O(n^{-1})$ de la coordenada física
de su índice: $(i-1/2)/n$, $(j-1/2)/n$, $k/n$ o $l/n$.
Esto vale también cuando las cardinalidades de las hojas difieren.

\paragraph{Núcleos límite entre familias distintas.}
Definimos $\mathcal K$ mediante los núcleos siguientes; los bloques en el orden inverso
son sus adjuntos y los bloques de una familia consigo misma son cero:
\[
\begin{array}{ll}
 RC:&\nu(r,c),\\
 RD:&\nu(r,r-d)\mathbf1_{0<r-d<1},\\
 RA:&\nu(r,a-r)\mathbf1_{0<a-r<1},\\
 CD:&\nu(c+d,c)\mathbf1_{0<c+d<1},\\
 CA:&\nu(a-c,c)\mathbf1_{0<a-c<1}.
\end{array}                                             \tag{D3.4}
\]
Estos cinco núcleos son independientes de ambas hojas.
En cambio,
\[
 \begin{gathered}
 \mathcal K_{DA}((d,p),(a,q))
 =\mathbf1_{p\ne q}\,
   \nu\left(\frac{a+d}{2},\frac{a-d}{2}\right)\mathbf1_\Omega(d,a),\\
 \Omega=\left\{(d,a):0<\frac{a+d}{2}<1,\ 0<\frac{a-d}{2}<1\right\}.
 \end{gathered}                                         \tag{D3.5}
\]
En efecto, $k=i-j$, $l=i+j-1$, y $k+l=2i-1$ es impar.
Esta condición no depende de $n\bmod2$.
La notación de $\Omega$ exige que ambos argumentos estén entre cero y uno.

\begin{proposicion}[Convergencia Hilbert--Schmidt]
Sea $f_0$ igual a $1/\sqrt2$ en ambas hojas R, a $-1/\sqrt2$ en
ambas hojas C y a cero en las familias D y A; $\|f_0\|=1$.
Con $P_0=|f_0\rangle\langle f_0|$, se tiene
\[
 \widehat J_n:=E_nJ_nE_n^*\longrightarrow
             J:=\mathcal K+P_0\quad\text{en Hilbert--Schmidt}.       \tag{D3.6}
\]
\end{proposicion}
\begin{proof}
El núcleo de una entrada de intersección es
$g_{ij}/\sqrt{\mu_{F,p}\mu_{G,q}}=ng_{ij}(1+O(n^{-1}))
=\nu(x_i,x_j)+O(n^{-1})$.
Fuera de bandas de anchura $O(n^{-1})$ alrededor de los segmentos, en número finito,
que delimitan las intersecciones, los indicadores coinciden
y los núcleos convergen uniformemente porque el perfil es Lipschitz.
Dentro de las bandas los núcleos permanecen acotados, y el área de las bandas tiende a cero.
Así cada bloque converge en $L^2$. La paridad de (D3.5) se conserva exacta.
La diagonal $\delta$ desaparece en norma Hilbert--Schmidt.
Finalmente $f_n=E_nv_n/\sqrt{2n}\to f_0$ en $\mathcal H$:
su valor en una hoja R es $\sqrt{m_{R,p}/n}\to1/\sqrt2$,
y análogamente en C. Por tanto sus proyectores convergen en esa norma.
\end{proof}

\paragraph{Lectura equivalente de la paridad.}
El cambio $s=(f_{F,0}+f_{F,1})/2$, $o=(f_{F,0}-f_{F,1})/2$
es unitario hacia $\bigoplus_F(L^2(I_F,dt)\oplus L^2(I_F,dt))$.
El sector simétrico tiene cuatro familias, y el bloque entre D y A es
$T(d,a)=\tfrac12\nu((a+d)/2,(a-d)/2)\mathbf1_\Omega$.
En el sector antisimétrico sólo queda el bloque
$\left(\begin{smallmatrix}0&-T\\-T^*&0\end{smallmatrix}\right)$ en las familias D y A.
Si se promediase la paridad antes de construir el límite, se perdería ese bloque.
